\section{On the role of the axioms}\label{sec:triviality}

\subsection{Structure of the result}

The proof of Theorem~\ref{thm:A} has three steps.

\begin{enumerate}[label=(\Alph*),leftmargin=*,itemsep=2pt]
\item \emph{Implied pair.} Axiom \axP{} defines $\Pi$; \axN{} and \axCM{}, applied with sure losses, define $\rho$ and
  give the Möbius form of the cash lines (Lemmas~\ref{lem:premium} and~\ref{lem:implied}).
\item \emph{Transfer.} Given \axP{} and \axN, \axCM{} holds if and only if $\rho$ and $\Pi$ are comonotonic
  additive (Lemma~\ref{lem:comon}), and \axLI{} passes from $\LR$ to $\rho$ and $\Pi$ (Lemma~\ref{lem:LI}).
\item \emph{Representation.} A law-invariant, comonotonic additive, uniformly continuous functional on an
  atomless $L^\infty$ is a signed Choquet integral \citep[Thm.~1]{www2020}; under monotonicity it is a
  distortion premium \citep{schmeidler1986,wyp1997}.
\end{enumerate}

Step (C) is classical; we reprove it in Section~\ref{sec:proofs} to keep the constants explicit and to show that
no continuity is needed in the monotone case. Steps (A) and (B) are specific to loss ratios. Together, axioms
\axP, \axN{} and \axCM{} state that $\LR$ is the ratio of a pair of comonotonic additive functionals that coincide
with the identity on constants, and that this pair can be read off $\LR$.

\subsection{Assumed and derived properties}

Table~\ref{tab:derived} lists the properties of the representation and their source. Axiom \axP{} is ordinal and
serves as a definition, \axN{} is a normalization, and \axLI, \axRC, \axRM{} select the class of representing
functions. The cardinal content of the axioms lies in \axCM: given \axP{} and \axN, it is equivalent to
comonotonic additivity of the implied premium and risk measure, which is the axiom that produces Choquet
representations in \citet{yaari1987}, \citet{schmeidler1986} and \citet{wyp1997}. The Choquet form of the numerator
and the denominator is therefore a consequence of comonotonic additivity, as in the single-functional case.

\begin{table}[ht]
\centering\small
\begin{tabularx}{\textwidth}{@{}>{\raggedright\arraybackslash}p{4.8cm}>{\raggedright\arraybackslash}X>{\raggedright\arraybackslash}p{2.0cm}@{}}
\toprule
\textbf{Property} & \textbf{Source} & \textbf{Status}\\
\midrule
$\Pi$ exists and is cash additive & \axP{} & definition\\
$\Pi(c)=c$, $\rho(c)=c$ & \axN{} and \axP & derived\\
$\rho$ defined on $L^\infty$; Möbius cash lines & \axCM{} with sure losses (Prop.~\ref{prop:readings}) & derived\\
$\rho,\Pi$ comonotonic additive & \axCM{} (Lemma~\ref{lem:comon}) & equivalent\\
$\rho,\Pi$ law invariant & \axLI{} (Lemma~\ref{lem:LI}) & assumed\\
$\rho,\Pi$ continuous or monotone & \axRC, \axRM{} & assumed\\
$h,g$ of bounded variation & Lemma~\ref{lem:bv} & derived\\
$\rho=I_h$, $\Pi=I_g$ & Lemmas~\ref{lem:layer} and~\ref{lem:rep} & derived\\
positive homogeneity, scale invariance & Lemma~\ref{lem:cashlip} & derived\\
association & Corollary~\ref{cor:assoc} & derived\\
uniqueness of $(h,g)$ & Proposition~\ref{prop:unique} & derived\\
\bottomrule
\end{tabularx}
\caption{Assumed and derived properties in Theorem~\ref{thm:A}.}\label{tab:derived}
\end{table}

\subsection{Relation to the single-functional theory}

The characterizations in the literature concern premiums and risk measures; they cannot be applied to $\LR$
directly, since a loss ratio is not additive and \axCM{} is a weighted-mean identity. What is specific to the
loss ratio is the construction of the implied pair, the derivation of additivity on all comonotone pairs, in
particular on the break-even set where the weights of \axCM{} are not identified (Remark~\ref{rem:primitive}),
uniqueness of $(h,g)$ also when $\LR\equiv1$, the impossibility of a monotone loss ratio
(Proposition~\ref{prop:F2}), the structure of the associated family (Corollary~\ref{cor:assoc}), and the version
for nonnegative losses (Section~\ref{sec:actuarial}), in which the price point is not observed and must be
extrapolated.

\subsection{Loss ratios excluded and included}

Ratios with a proportional loading, $\E/((1+\theta)\E)$, violate \axN; scenario-weighted ratios $\E_Q/\E$ violate
\axLI; the ratio $\max(I_{\sqrt t},I_{\min(1,1.5t)})/\E$ violates \axCM, its numerator being cash additive but not
comonotonic additive; and ordinal transforms such as $(\rho/\Pi)^3$ satisfy \axP{} and \axN{} but violate \axCM.
The class characterized contains, among others, the expected loss over any distortion premium
(Example~\ref{ex:PH}), ratios of expected shortfalls (Example~\ref{ex:ES}), value-at-risk over expected shortfall,
and ratios that take values on both sides of $1$.
